\chapter{DESIGN AND IMPLEMENTATION}

\section{Design Changes}

% TODO: Add content for design changes section

\section{Implementation Challenges}

% TODO: Add content for implementation challenges section

\section{Formative User Testing Observations}

Formative user testing was conducted with Grade 6 students from Inchican Elementary School using a playable prototype that included the Tutorial Level, Bubble Shooter Level 1, and the Sorting Mini Game. Testing was conducted face-to-face, allowing direct observation of player behavior, engagement, and comprehension. Students demonstrated an overall understanding of gameplay mechanics and responded positively to visual cues and immediate feedback, despite their current science curriculum focusing on topics outside states of matter.

Feedback indicated that some levels were perceived as too easy and quickly memorized, suggesting a need for increased challenge variety. Students also commented on improving visual consistency in the sorting minigame and adding more objects to extend engagement. Language accessibility emerged as a key consideration, with requests for Filipino or Tagalog text. Overall, students expressed enjoyment of the game experience and interest in playing a completed version. These observations will inform future refinements to difficulty balancing, visual design, and language support.
